\documentclass[10pt,a4paper]{amsart}
\usepackage[margin=19mm]{geometry}
\usepackage{amsmath,amssymb,mathtools}
\usepackage[scr=rsfs,cal=euler]{mathalfa}
\usepackage{xfrac}
\usepackage[unicode,hidelinks,bookmarks=false]{hyperref}
\newcommand{\ie}{i.e.\ }
\newcommand{\cf}{cf.\ }
\newcommand{\resp}{respectively\ }
\newcommand{\Subsection}{\subsection}
\providecommand{\nobreakdash}{\nobreak\hbox{-}}
\setlength{\emergencystretch}{2em}
\multlinegap0pt
\usepackage{bm}
\usepackage{bbm}
\usepackage{dsfont}
\usepackage{xspace}
\usepackage{cancel}
\usepackage{mathtools}
\usepackage{amsthm}
\makeatletter
\@ifundefined{assumption}{\newtheorem{assumption}{Assumption}}{}
\@ifundefined{lem}{\newtheorem{lem}{Lemma}}{}
\@ifundefined{thm}{\newtheorem{thm}{Theorem}}{}
\makeatother
\title[Global Existence and Pathwise Uniqueness for a Stochastic Pa]{Global Existence and Pathwise Uniqueness for a Stochastic Parabolic-Parabolic Keller-Segel System}
\author{Jinhuan Wang, Qian Li, Hui Huang}
\date{Source: arXiv:2607.24128v1 (CC BY 4.0)}
\begin{document}
\maketitle

\noindent\emph{Source and licence.} Global Existence and Pathwise Uniqueness for a Stochastic Parabolic-Parabolic Keller-Segel System. Version arXiv:2607.24128v1 (CC BY 4.0), distributed under the Creative Commons Attribution 4.0 International licence, as recorded in the arXiv licence metadata for this exact version and shown on the abstract page https://arxiv.org/abs/2607.24128v1. Full rights notice: licenses/arxiv-2607.24128-CC-BY-4.0.txt.\par\smallskip

\noindent\textit{Editorial scope.} Counted unit: the Thm whose source label is \texttt{thm\_estimate} in the article (source0.tex lines 1174--1181, source label \texttt{thm\_estimate}), delivered together with the article's own complete original proof of it (source0.tex lines 1183--1520, one \texttt{proof} environment of 338 lines). No mathematics is written, repaired or re-derived: statement and proof are the article's own text line for line. Required context retained uncounted: equ1 (lines 131--145); ass (lines 251--265); equ-v\_t (lines 408--413); lem\_v (lines 414--446). Every definition, display and label of those lines is retained unchanged. Omitted from this package and not redistributed: 1--130, 146--250, 266--407, 447--1173, 1182--1182, 1521--1574. Nothing inside a retained excerpt is omitted or altered: every retained body is delivered exactly as the article's lines are written, so no omission receipt of any kind accompanies this package. Citations: the retained excerpts carry no citation command at all, so no bibliography entry of the article is transcribed and no reference is redistributed. Reference adaptation: none was needed and none was made; every reference inside the delivered unit resolves inside it, so no unresolved reference or citation is produced. Printed slips kept literal: no printed word, symbol, hypothesis or number of the counted statement or of its proof was altered. Layout and class: the article's own class is \texttt{amsart} with packages including geometry, enumerate, amssymb, amsmath, xcolor, lineno, hyperref, enumitem; the wrapper is the lane's pinned \texttt{amsart} editorial wrapper at 10pt on A4, which restates the theorem-like environments the retained text needs and adds the article's own macro definitions that the retained text needs (none), with extra wrapper packages (bm, bbm, dsfont, xspace, cancel, mathtools). The delivered numbering is the unit's own. Assets: no retained span contains a figure environment, a graphics-inclusion command, a PDF-page inclusion, a table or a TikZ picture; no image file is copied into this package, the package declares no asset dependency and no image file is redistributed with it. Licence: version arXiv:2607.24128v1 of this article is distributed under the Creative Commons Attribution 4.0 International licence, as recorded in the arXiv licence metadata for this exact version and shown on the abstract page https://arxiv.org/abs/2607.24128v1. A full rights notice is delivered at licenses/arxiv-2607.24128-CC-BY-4.0.txt. Characters: every retained excerpt is pure ASCII, so the delivered engine, XeLaTeX with its default Latin Modern text font, reports no missing character.
\begin{equation}\label{equ1}
\begin{cases}
du = \big(\Delta u-\nabla\cdot(u\nabla v)\big)\,dt+\sigma(u)\,dW(t),
& x\in\mathcal{O},\ t>0,\\
dv = \big(\Delta v-v+u\big)\,dt,
& x\in\mathcal{O},\ t>0,\\
\dfrac{\partial u}{\partial \nu}
=
\dfrac{\partial v}{\partial \nu}
=0,
& x\in\partial\mathcal{O},\ t>0,\\
u(x,0)=u_0(x),\quad v(x,0)=v_0(x),
& x\in\mathcal{O},
\end{cases}
\end{equation}

\begin{assumption}\label{ass}
Assume that
\[
2< p<\infty,
\qquad
1<r<\frac{p}{p-2}.
\]
Let $\sigma:H^1(\mathcal O)\to L_2(U,L^2(\mathcal O))$
be defined by
\[
(\sigma(u)e_k)(x):=b_k\|u\|_{L^p(\mathcal O)}^{r}u(x),
\qquad u\in H^1(\mathcal O),\ k\in\mathbb N,
\]
where $(b_k)_{k\ge1}\in \ell^2(\mathbb R).$
\end{assumption}

\begin{equation}\label{equ-v_t}
\begin{cases}
v_t=\Delta v-v+g, & x\in\mathcal {O},\ t>0, \\\frac{\partial v}{\partial \nu} = 0,&x\in\partial\mathcal{O},t> 0,\\
v(x,0)=v_0(x), & x\in\mathcal {O}.
\end{cases}
\end{equation}

\begin{lem}\label{lem_v}
	Let $\mathcal O\subset\mathbb R^2$ be a bounded smooth domain and let $v$ be the mild solution of \eqref{equ-v_t}.
    	\begin{itemize}
		\item[(i)] Assume that $v_0\in W^{1,\infty}(\mathcal O)$ and
		$g\in L^2(0,T;H^1(\mathcal O))$. Then there exists a constant $C>0$, independent of $T$, such that for all $t\in[0,T]$,
\begin{equation}\label{v_infty}
			\int_0^t \|\nabla v(s)\|_{L^\infty(\mathcal O)}^2\,ds
			\le
			C\left(
			t\|\nabla v_0\|_{L^\infty(\mathcal O)}^2
			+(t+t^2)
			\int_0^t \|g(s)\|_{H^1(\mathcal O)}^2\,ds
			\right).
		\end{equation}
		\item[(ii)] Assume that
		\[
		1\le m\le q\le \infty,\qquad 1\le \gamma<\infty,
		\qquad \frac1q>\frac1m-\frac12,
		\]
		and suppose that $v_0\in W^{1,q}(\mathcal O), \,g\in L^\gamma(0,T;L^m(\mathcal O)).$
		Then, for all $t\in[0,T]$, there exists a positive constant $C$,
		independent of $T$, such that
	\begin{equation}\label{v_r,q}
			\int_0^t \|\nabla v(s)\|_{L^q(\mathcal O)}^\gamma \,ds
			\le
			C\left(
			t\|\nabla v_0\|_{L^q(\mathcal O)}^\gamma
			+
			\int_0^t \|g(s)\|_{L^m(\mathcal O)}^\gamma \,ds
			\right).
		\end{equation}
	\end{itemize}
\end{lem}

  \begin{thm}\label{thm_estimate}
      Let $(u,\tau_N)$ be a local strong solution  of  system (\ref{equ1}). Under Assumption \ref{ass},  for every $T>0$ and every $\alpha\in\big(\max\{0,\frac{2-r}{2}\},\frac{1}{2}\big)$, there exists a constant $C(T)>0$, independent of $N$, such that 
      \begin{align}\label{est_C(T)}
	    \mathbb E\Big(\sup_{0\le t\le T}(1+\|u(t\wedge\tau_N)\|_{L^2(\mathcal{O})}^2)^\alpha\Big)+\mathbb{E}\Big(\int_0^{T\wedge{\tau_N}}
		\frac{\|\nabla u(s)\|_{L^2(\mathcal{O})}^2}{(1+\|u(s)\|_{L^2(\mathcal{O})}^2)^{1-\alpha}}\,ds\Big)
	\le C(T).
	\end{align}
  \end{thm} 

  \begin{proof}
    	 Applying It\^o's formula to $\|u(t\wedge{\tau_N})\|_{L^2(\mathcal O)}^2$, for convenience, we denote the $L^2(\mathcal{O})$ norm of $u$ and $\nabla u$  simply by $\|u\|_2$ and $\|\nabla u\|_2$ respectively and we obtain
        \begin{align}\label{es_L2}
		\nonumber&\|u(t\wedge{\tau_N})\|_2^2-\|u_0\|_2^2\\
		=&2\int_0^{t\wedge{\tau_N}}\big(u(s),\Delta u(s)\big)_2\,ds
		-2\int_0^{t\wedge{\tau_N}}\big(u(s),\nabla\cdot(u(s)\nabla v(s))\big)_2\,ds
		\nonumber\\
&+B_0^2\int_0^{t\wedge{\tau_N}}\|u(s)\|_p^{2r}\|u(s)\|_2^2\,ds+2\sum_{k=1}^\infty\int_0^{t\wedge{\tau_N}}b_k\|u(s)\|_p^r\|u(s)\|_2^2dW_k(s).
	\end{align}
  We introduce the Lyapunov function $\Phi(x)=(1+x)^\alpha$ with $\alpha\in\big(\max\{0,\frac{2-r}{2}\},\frac{1}{2}\big).$ Applying It\^o's formula to $\Phi(\|u(t\wedge{\tau_N})\|_2^2)$ and combining \eqref{es_L2} yields
	\begin{align*}
		\nonumber&\big(1+\|u(t\wedge{\tau_N})\|_2^2\big)^\alpha
+2\alpha\int_0^{t\wedge{\tau_N}}\frac{\|\nabla u(s)\|_2^2}{(1+\|u(s)\|_2^2)^{1-\alpha}}\,ds
		\\\nonumber
		=&(1+\|u_0\|_2^2)^\alpha
		+2\alpha\int_0^{t\wedge{\tau_N}}
		\frac{\int_{\mathcal O}u(s)\nabla u(s)\cdot \nabla v(s)dx}
		{(1+\|u(s)\|_2^2)^{1-\alpha}}\,ds
+\alpha\int_0^{t\wedge{\tau_N}}
		\frac{B_0^2\|u(s)\|_p^{2r}\|u(s)\|_2^2}
		{(1+\|u(s)\|_2^2)^{1-\alpha}}\,ds
		\\
		&-2\alpha(1-\alpha)\int_0^{t\wedge{\tau_N}}
		\frac{{B_0^2}\|u(s)\|_p^{2r}\|u(s)\|_2^4}
		{(1+\|u(s)\|_2^2)^{2-\alpha}}\,ds
+2\alpha\sum_{k=1}^\infty\int_0^{t\wedge{\tau_N}}
		\frac{b_k\|u(s)\|_p^r\|u(s)\|_2^2}
		{(1+\|u(s)\|_2^2)^{1-\alpha}}dW_k(s).
	\end{align*}
    Taking $2<m\le \frac{(r+2)p}{p+r}$ and  by H\"older's and Young's inequalities,  we get  for any $\varepsilon_1>0$,
	\begin{align}\label{ine_v}
	    2\alpha\int_{\mathcal O}u\nabla u\cdot \nabla v\,dx
	\le
	\varepsilon_1\|\nabla u\|_2^2
	+C\|u\|_m^2\|\nabla v\|_{\frac{2m}{m-2}}^2.
	\end{align}
	Substituting \eqref{ine_v} into the above equality, we obtain
	\begin{align}
		\label{L3}
		\nonumber&(1+\|u(t\wedge{\tau_N})\|_2^2)^\alpha
		+(2\alpha-\varepsilon_1)\int_0^{t\wedge{\tau_N}}\frac{\|\nabla u(s)\|_2^2}{(1+\|u(s)\|_2^2)^{1-\alpha}}\,ds\\\nonumber&+(\alpha-2\alpha^2)\int_0^{t\wedge{\tau_N}}
		\frac{{B_0^2}\|u(s)\|_p^{2r}\|u(s)\|_2^4}
		{(1+\|u(s)\|_2^2)^{2-\alpha}}\,ds
		\nonumber\\
		\le&(1+\|u_0\|_2^2)^\alpha
		+C\int_0^{t\wedge{\tau_N}}
		\frac{\|u(s)\|_m^2\|\nabla v(s)\|_{\frac{2m}{m-2}}^2}
		{(1+\|u(s)\|_2^2)^{1-\alpha}}\,ds
		+\alpha\int_0^{t\wedge{\tau_N}}
		\frac{{B_0^2}\|u(s)\|_p^{2r}\|u(s)\|_2^2}
		{(1+\|u(s)\|_2^2)^{2-\alpha}}\,ds
		\nonumber\\
		&
+2\alpha\sum_{k=1}^\infty\int_0^{t\wedge{\tau_N}}
		\frac{b_k\|u(s)\|_p^r\|u(s)\|_2^2}
		{(1+\|u(s)\|_2^2)^{1-\alpha}}dW_k(s).
	\end{align}
    Next, we use the Gagliardo-Nirenberg inequality, the interpolation inequality, and Young's inequality to estimate the  integral terms on the right-hand side in \eqref{L3} so that they can be absorbed into the left-hand side in \eqref{L3}. 
    
    We claim that for any $\varepsilon_2,\varepsilon_3>0$,
	\begin{equation}
		\label{es_p,2}
		\alpha{B_0^2}\|u\|_p^{2r}\|u\|_2^2
		\le
		\varepsilon_2\|\nabla u\|_2^2\|u\|_2^2
		+\varepsilon_3 \alpha{B_0^2}\|u\|_p^{2r}\|u\|_2^4
		+C(\varepsilon_2,\varepsilon_3,\alpha,{B_0^2}).
	\end{equation}
	Indeed,   if $\|u\|_2^2\ge \varepsilon_3^{-1}$, then
	\[
	\alpha{B_0^2}\|u\|_p^{2r}\|u\|_2^2
	\le
	\varepsilon_3 \alpha{B_0^2}\|u\|_p^{2r}\|u\|_2^4.
	\]
    If $\|u\|_2^2<\varepsilon_3^{-1}$, by the Gagliardo--Nirenberg inequality in two dimensions,
	\begin{align*}
		\|u\|_p^{2r}\|u\|_2^2
		&\le
		C\|\nabla u\|_2^{\frac{2r(p-2)}{p}}\|u\|_2^{\frac{2p+4r}{p}}
		+C\|u\|_2^{2r+2}\\
		&=
		C\|u\|_2^{\frac{2(p+4r-rp)}{p}}
		\big(\|\nabla u\|_2^2\|u\|_2^2\big)^{\frac{r(p-2)}{p}}
		+C\|u\|_2^{2r+2}.
	\end{align*}
	Since $p>2$ and $r<\frac{p}{p-2}$, we have
	\[
	\frac{r(p-2)}{p}<1,
	\qquad
	\frac{2(p+4r-rp)}{p}>0.
	\]
	 Therefore,
	\[
	\|u\|_p^{2r}\|u\|_2^2
	\le
	C(\varepsilon_3)\big(\|\nabla u\|_2^2\|u\|_2^2\big)^{\frac{r(p-2)}{p}}
	+C(\varepsilon_3),
	\]
	and Young's inequality yields
	\[
	\alpha{B_0^2}\|u\|_p^{2r}\|u\|_2^2
	\le
	\varepsilon_2\|\nabla u\|_2^2\|u\|_2^2
	+C(\varepsilon_2,\varepsilon_3,\alpha,{B_0^2}).
	\]
	Combining the above two cases, we obtain \eqref{es_p,2}.
	Furthermore, using
	\[
	\frac{\|u\|_2^2}{(1+\|u\|_2^2)^{2-\alpha}}
	\le
	\frac{1}{(1+\|u\|_2^2)^{1-\alpha}},
	\qquad
	\frac{1}{(1+\|u\|_2^2)^{2-\alpha}}\le 1,
	\]
	we infer that
	\begin{align}
		\label{es_a,q,2}
		\alpha\frac{{B_0^2}\|u\|_p^{2r}\|u\|_2^2}{(1+\|u\|_2^2)^{2-\alpha}}
		\le
		\varepsilon_2\frac{\|\nabla u\|_2^2}{(1+\|u\|_2^2)^{1-\alpha}}
		+\varepsilon_3 \alpha\frac{{B_0^2}\|u\|_p^{2r}\|u\|_2^4}{(1+\|u\|_2^2)^{2-\alpha}}
+C(\varepsilon_2,\varepsilon_3).
	\end{align}
	We next estimate the term involving $\|u\|_m^2\|\nabla v\|_{\frac{2m}{m-2}}^2$.
	By Young's inequality, we have
	\begin{align}
		\label{u-v-term}
		C\frac{\|u\|_m^2\|\nabla v\|_{\frac{2m}{m-2}}^2}{(1+\|u\|_2^2)^{1-\alpha}}
		\le
		C\frac{\|u\|_m^{\frac{2(2-\alpha)}{1-\alpha}}}{(1+\|u\|_2^2)^{2-\alpha}}
		+\|\nabla v\|_{\frac{2m}{m-2}}^{2(2-\alpha)}.
	\end{align}
	Since $2<m\le \frac{(r+2)p}{p+r},$
	taking $\theta=\frac{r}{r+2}$ yields $\frac1m\ge \frac{p+r}{(r+2)p}= \frac{\theta}{p}+\frac{1-\theta}{2}.$
	Hence, by interpolation inequality on bounded domains, it holds 
	\begin{align}\label{interpolation_1}
	    \|u\|_m^{\frac{2(2-\alpha)}{1-\alpha}}\le C\|u\|_{\frac{(r+2)p}{p+r}}^{\frac{2(2-\alpha)}{1-\alpha}}
	\le
	C\|u\|_p^{\frac{2(2-\alpha)r}{(1-\alpha)(r+2)}}
	\|u\|_2^{\frac{4(2-\alpha)}{(1-\alpha)(r+2)}}
	=
	C\big(\|u\|_p^{2r}\|u\|_2^4\big)^{\frac{2-\alpha}{(1-\alpha)(r+2)}}.
	\end{align}
	Since $\alpha<\frac{1}{2}<\frac{r}{r+1}$, we have $\frac{2-\alpha}{(1-\alpha)(r+2)}<1.$
	Therefore, using Young's inequality arrives
	\begin{align}
		\label{u_m_a}
		C\frac{\|u\|_m^{\frac{2(2-\alpha)}{1-\alpha}}}{(1+\|u\|_2^2)^{2-\alpha}}
		\le
		\varepsilon_4 \alpha\frac{{B_0^2}\|u\|_p^{2r}\|u\|_2^4}{(1+\|u\|_2^2)^{2-\alpha}}
		+C(\varepsilon_4).
	\end{align}
	For $\int_0^{t\wedge{\tau_N}}\|\nabla v(s)\|_{\frac{2m}{m-2}}^{2(2-\alpha)}\,ds$,  by Lemma \ref{lem_v}, we know 
	\begin{equation}
		\int_0^{t\wedge{\tau_N}}\|\nabla v(s)\|_{\frac{2m}{m-2}}^{2(2-\alpha)}\,ds
		\le
		Ct\|\nabla v_0\|_{\frac{2m}{m-2}}^{2(2-\alpha)}+C\int_0^{t\wedge{\tau_N}}\|u(s)\|_m^{2(2-\alpha)}\,ds.
	\end{equation}
	We further claim that for any $\varepsilon_5>0$,
	\begin{equation}
		\label{ine_u,m}
C\|u\|_m^{2(2-\alpha)}
\le\varepsilon_5
		\frac{\alpha{B_0^2}\|u\|_p^{2r}\|u\|_2^4}{(1+\|u\|_2^2)^{2-\alpha}}
		+
		C(\varepsilon_5).
	\end{equation}
    Indeed, if  $\|u\|_2^2\le 1$,
	then
	\[
	(1+\|u\|_2^2)^{2-\alpha}\le 2^{2-\alpha},
	\]
	and hence
	\[
	\frac{\alpha{B_0^2}\|u\|_p^{2r}\|u\|_2^4}{(1+\|u\|_2^2)^{2-\alpha}}
	\ge
	2^{-(2-\alpha)}\alpha{B_0^2}\|u\|_p^{2r}\|u\|_2^4.
	\]
	Similarly with \eqref{interpolation_1}, we use again the interpolation inequality with $\theta=\frac{r}{r+2}$ to get
	\[
	\|u\|_m^{2(2-\alpha)}
	\le C\|u\|_{\frac{(r+2)p}{p+r}}^{2(2-\alpha)}\le
	C\|u\|_p^{\frac{2(2-\alpha)r}{r+2}}
	\|u\|_2^{\frac{4(2-\alpha)}{r+2}}
	=
	C(\|u\|_p^{2r}\|u\|_2^4)^{\frac{2-\alpha}{r+2}}.
	\]
	Since $\frac{2-\alpha}{r+2}<1$, Young's inequality implies
	\begin{align}
	    \|u\|_m^{2(2-\alpha)}
	\le
	\varepsilon_5\frac{\alpha{B_0^2}\|u\|_p^{2r}\|u\|_2^4}{(1+\|u\|_2^2)^{2-\alpha}}
	+
	C(\varepsilon_5).
	\end{align}
	 If $\|u\|_2^2>1$, then
	\[
	(1+\|u\|_2^2)^{2-\alpha}\le 2^{2-\alpha}\|u\|_2^{2(2-\alpha)},
	\]
	and therefore
	\[
	\frac{\alpha{B_0^2}\|u\|_p^{2r}\|u\|_2^4}{(1+\|u\|_2^2)^{2-\alpha}}
	\ge
	2^{-(2-\alpha)}\alpha{B_0^2}\|u\|_p^{2r}\|u\|_2^{2\alpha}.
	\]
	Since
	\[
	m\le \frac{(r+2)p}{p+r}<\frac{2p(r+\alpha)}{2r+\alpha p},
	\]
	we may take $\theta=\frac{r}{r+\alpha}$ so that
	\[
	\frac1m\ge \frac{\theta}{p}+\frac{1-\theta}{2},
	\]
	Consequently, using the interpolation inequality yields
	\[
	\|u\|_m^{2(2-\alpha)}
	\le
	C\|u\|_p^{\frac{2(2-\alpha)r}{r+\alpha}}
	\|u\|_2^{\frac{2\alpha(2-\alpha)}{r+\alpha}}
	=
	C(\|u\|_p^{2r}\|u\|_2^{2\alpha})^{\frac{2-\alpha}{r+\alpha}}.
	\]
	Since $\frac{2-r}{2}<\alpha$, we have $\frac{2-\alpha}{r+\alpha}<1.$
	Thus, Young's inequality  gives
	\begin{align}
	    \|u\|_m^{2(2-\alpha)}
	\le
	\varepsilon_5\frac{\alpha{B_0^2}\|u\|_p^{2r}\|u\|_2^4}{(1+\|u\|_2^2)^{2-\alpha}}
	+
	C(\varepsilon_5).
	\end{align}
	Combining the above two cases, \eqref{ine_u,m} holds.
Thus, it follows  that
\begin{equation}
		\label{ine_v,m}
		\int_0^{t\wedge{\tau_N}}\|\nabla v(s)\|_{\frac{2m}{m-2}}^{2(2-\alpha)}\,ds
		\le
\varepsilon_5\int_0^{t\wedge{\tau_N}}\frac{\alpha{B_0^2}\|u(s)\|_p^{2r}\|u(s)\|_2^4}{(1+\|u(s)\|_2^2)^{2-\alpha}}\,ds+C(v_0,\varepsilon_5,t).
	\end{equation}
	Substituting  \eqref{es_a,q,2}, \eqref{u-v-term}, \eqref{u_m_a} and  \eqref{ine_v,m}
	into \eqref{L3}, we arrive at
	\begin{align}
		\label{La}
		\nonumber&\big(1+\|u(t\wedge{\tau_N})\|_2^2\big)^\alpha
		+(2\alpha-\varepsilon_1-\varepsilon_2)\int_0^{t\wedge{\tau_N}}
		\frac{\|\nabla u(s)\|_2^2}{(1+\|u(s)\|_2^2)^{1-\alpha}}\,ds\\\nonumber&+\bigl(\alpha-2\alpha^2-\varepsilon_3-\varepsilon_4-\varepsilon_5\bigr)
		\int_0^{t\wedge{\tau_N}}
		\frac{{B_0^2}\|u(s)\|_p^{2r}\|u(s)\|_2^4}
		{(1+\|u(s)\|_2^2)^{2-\alpha}}\,ds
		\nonumber\\
		\le&\big(1+\|u_0\|_2^2\big)^\alpha+2\alpha\sum_{k=1}^\infty\int_0^{t\wedge{\tau_N}}
		\frac{b_k\|u(s)\|_p^r\|u(s)\|_2^2}{(1+\|u(s)\|_2^2)^{1-\alpha}}dW_k(s)
		+C(\varepsilon_1, \varepsilon_2,\varepsilon_3,\varepsilon_4,\varepsilon_5,t,v_0).
	\end{align}
    Choosing $\varepsilon_i>0$, $1\le i\le 5$ sufficiently small such that
	\[
	2\alpha-\varepsilon_1-\varepsilon_2>0,
	\qquad
	\alpha-2\alpha^2-\varepsilon_3-\varepsilon_4-\varepsilon_5>0.
	\]
	Taking expectations in \eqref{La}, and using the definition of $\tau_N$, we see that the stochastic integral is a square-integrable  martingale with zero expectation. Then, we obtain
	\begin{align}
		\label{exp_u}
		\mathbb E\Big(\int_0^{t\wedge{\tau_N}}
		\frac{{B_0^2}\|u(s)\|_p^{2r}\|u(s)\|_2^4}
		{(1+\|u(s)\|_2^2)^{2-\alpha}}\,ds\Big)
		\le C(t,v_0,u_0).
	\end{align}
     Furthermore, taking the supremum over $t\in[0,T]$ in \eqref{La}, then taking expectations, we get
	\begin{align}
		\label{est_sup_u}
		&\mathbb E\Big(\sup_{0\le t\le T}(1+\|u(t\wedge{\tau_N})\|_2^2)^\alpha\Big)+\mathbb{E}\Big(\int_0^{T\wedge{\tau_N}}
		\frac{\|\nabla u(s)\|_2^2}{(1+\|u(s)\|_2^2)^{1-\alpha}}\,ds\Big)
		\nonumber\\
		&\le
		\mathbb E\big(\big(1+\|u_0\|_2^2\big)^\alpha\big)
		+2\alpha\mathbb E\Big(\sup_{0\le t\le T}
		\Big|
		\sum_{k=1}^\infty\int_0^{t\wedge{\tau_N}}
		\frac{b_k\|u(s)\|_p^r\|u(s)\|_2^2}{(1+\|u(s)\|_2^2)^{1-\alpha}}dW_k(s)
		\Big|\Big)
		\nonumber\\
&+C(T,v_0,u_0).
	\end{align}
	By the Burkholder-Davis-Gundy inequality and Young's inequality,  we have 
	\begin{align}
		\label{BDG-estimate}
		&2\alpha\mathbb E\Big(\sup_{0\le t\le T}
		\Big|
		\sum_{k=1}^\infty\int_0^{t\wedge{\tau_N}}
		\frac{b_k\|u(s)\|_p^r\|u(s)\|_2^2}{(1+\|u(s)\|_2^2)^{1-\alpha}}dW_k(s)
		\Big|\Big)
		\nonumber\\
		\le&
		C\mathbb E\Big[\Big(
		\int_0^{T\wedge{\tau_N}}
		\frac{{B_0^2}\|u(s)\|_p^{2r}\|u(s)\|_2^4}
		{(1+\|u(s)\|_2^2)^{2-2\alpha}}\,ds
		\Big)^{1/2}\Big]
		\nonumber\\
		=&
		C\mathbb E\Big[\Big(
		\int_0^{T\wedge{\tau_N}}
		\big(1+\|u(s)\|_2^2\big)^\alpha
		\frac{{B_0^2}\|u(s)\|_p^{2r}\|u(s)\|_2^4}
		{(1+\|u(s)\|_2^2)^{2-\alpha}}\,ds
		\Big)^{1/2}\Big]
		\nonumber\\
		\le&
		C\mathbb E\Big[
		\Big(
		\sup_{0\le t\le T}(1+\|u(t\wedge{\tau_N})\|_2^2)^\alpha
		\Big)^{1/2}
		\Big(
		\int_0^{T\wedge{\tau_N}}
		\frac{{B_0^2}\|u(s)\|_p^{2r}\|u(s)\|_2^4}
		{(1+\|u(s)\|_2^2)^{2-\alpha}}\,ds
		\Big)^{1/2}
		\Big]
		\nonumber\\
		\le&
		\frac12\mathbb E\Big(\sup_{0\le t\le T}(1+\|u(t\wedge{\tau_N})\|_2^2)^\alpha\Big)
		+C\mathbb E\Big(\int_0^{T\wedge{\tau_N}}
		\frac{{B_0^2}\|u(s)\|_p^{2r}\|u(s)\|_2^4}
		{(1+\|u(s)\|_2^2)^{2-\alpha}}\,ds\Big)
		\nonumber\\
		\le&
		\frac12\mathbb E\Big(\sup_{0\le t\le T}(1+\|u(t\wedge{\tau_N})\|_2^2)^\alpha\Big)
		+C(T,v_0,u_0),
	\end{align}
	where we have used the estimate \eqref{exp_u} in the last step.  Substituting \eqref{BDG-estimate} into \eqref{est_sup_u}, we conclude that
	\begin{align}
	    \mathbb E\Big(\sup_{0\le t\le T}(1+\|u(t\wedge\tau_N)\|_2^2)^\alpha\Big)+\mathbb{E}\Big(\int_0^{T\wedge{\tau_N}}
		\frac{\|\nabla u(s)\|_2^2}{(1+\|u(s)\|_2^2)^{1-\alpha}}\,ds\Big)
	\le C(T,v_0,u_0).
	\end{align}
    This completes the proof of Theorem \ref{thm_estimate}.
     \end{proof}

\end{document}
